\documentclass[10pt,letterpaper]{article}

\usepackage[spanish,es-noquoting]{babel}
\usepackage[letterpaper,top=1.45cm,bottom=1.55cm,left=1.85cm,right=1.85cm]{geometry}
\usepackage{tabularx}
\usepackage{array}
\usepackage{setspace}
\usepackage{microtype}
\usepackage[hidelinks]{hyperref}

\setlength{\parindent}{0pt}
\setlength{\parskip}{0.38em}
\renewcommand{\arraystretch}{1.03}
\pagestyle{empty}

\newcommand{\blankline}[1]{\makebox[#1]{\hrulefill}}
\newcommand{\field}[2]{#1: \blankline{#2}}

\begin{document}

{\large\bfseries Carta de apoyo institucional}\\[-0.25em]
{\small Concurso Innovación Docente AIEP 2026}

\vspace{0.12em}
\hrule height 0.7pt
\vspace{0.45em}

\begin{tabularx}{\textwidth}{>{\bfseries}p{3.45cm}X}
Sede & Osorno \\
Escuela & Ingeniería, Energía y Tecnología \\
Proyecto & Aula Subtitulada AIEP: subtítulos flotantes en tiempo real para clases inclusivas \\
Línea & Innovación en el Aula \\
Postulante & Diego Matías Obando Aguilera \\
\end{tabularx}

\vspace{0.35em}

Osorno, \blankline{2.0cm} de \blankline{3.0cm} de 2026

\vspace{0.2em}

Señores/as\\
\textbf{Comité Concurso Innovación Docente 2026}\\
AIEP\\
Presente

\vspace{0.2em}

De mi consideración:

Por medio de la presente, manifiesto mi apoyo institucional a la postulación del proyecto \textbf{``Aula Subtitulada AIEP: subtítulos flotantes en tiempo real para clases inclusivas''}, liderado por el docente \textbf{Diego Matías Obando Aguilera}, en el marco del \textbf{Concurso Innovación Docente AIEP 2026}, línea \textbf{Innovación en el Aula}.

La iniciativa propone implementar y validar en aula una solución de accesibilidad que permite transformar la voz del docente en subtítulos visibles en tiempo real, mediante una aplicación de apoyo pedagógico pensada para reducir barreras de comprensión, participación y acceso a instrucciones durante la clase. El proyecto se alinea con el Diseño Universal para el Aprendizaje, al ofrecer una vía visual complementaria para todo el curso.

El origen de esta propuesta surge a partir de una oportunidad identificada por \textbf{Elías Alberto Silva Carrasco}, Jefe de Escuela de Ingeniería, Energía y Tecnología de sede Osorno, quien planteó avanzar hacia una herramienta que apoyara a estudiantes con barreras auditivas mediante subtitulado o traducción de lo hablado por docentes. A partir de esa orientación, Diego Matías Obando Aguilera lideró el desarrollo del prototipo funcional y propone su implementación pedagógica, evaluación y mejora durante el segundo semestre de 2026.

Como Dirección, valoramos esta postulación por su pertinencia institucional, su foco inclusivo, su potencial de escalamiento a otras asignaturas y sedes, y su contribución a metodologías de enseñanza-aprendizaje apoyadas por tecnología. Asimismo, comprometemos disposición para facilitar la coordinación académica necesaria, definir la sección de aplicación, resguardar la implementación y apoyar la socialización de resultados en la sede.

Sin otro particular, y esperando una favorable acogida, saluda atentamente,

\vfill

\begin{center}
\begin{tabular}{>{\bfseries}p{1.8cm}p{7.2cm}}
Firma: & \rule{0pt}{1.2cm}\hrulefill \\[0.15em]
Nombre: & \hrulefill \\[0.15em]
Cargo: & \hrulefill \\
\end{tabular}
\end{center}

\end{document}
